% !PRINTED 478
\weberpage{478}

To find the characteristic properties of the Galois group, we have only to express the theorems \S.\,146, 1. to 4. somewhat differently. Of a function of the $m$ roots which remains unchanged under any permutation we say that it \gesperrt{admits this permutation}, and thus obtain:

a) \gesperrt{Every rational equation in $\Omega$ that holds between the $m$ roots of $F(x)$ remains valid when the roots are subjected to any permutation of the Galois group.}

b) \gesperrt{Every rational function in $\Omega$ of the $m$ roots of $F(x)$ that admits all the permutations of the Galois group is a number in $\Omega$.}

For if we express the roots of $F(x)$ as rational functions of $\varrho$, then a function of the roots $\alpha$ passes into a function $\varphi(\varrho)$. If a permutation of the Galois group has first been carried out, one obtains one of the conjugate numbers $\varphi(\varrho_a)$, which arises from $\varphi(\varrho)$ by a substitution $\sigma_a = (\varrho, \varrho_a)$. Now if $\varphi(\varrho) = 0$, then by the theorem \S.\,146, 1. all the conjugates $\varphi(\varrho_a) = 0$ as well, whereby a) is proved; and if the conjugate quantities $\varphi(\varrho_a)$ are all equal to one another, then their common value is, by \S.\,146, 4., contained in $\Omega$, whereby b) is proved.

To a) and b) a third statement is added:

c) \gesperrt{If any permutation $\pi$ of the roots of $F(x)$ is applicable to all rational equations in $\Omega$ that hold between the roots, then $\pi$ belongs to the Galois group, and the Galois group can therefore also be defined as the totality of all permutations that are applicable to all the rational equations between the roots.}

For by \S.\,143 one can take $\varrho$ to be a rational (e.\,g. linear) function of the $m$ quantities $\alpha, \alpha_1, \alpha_2 \ldots \alpha_{m-1}$ in such a way that all the $\Pi(m)$ values which one obtains from it by the $\Pi(m)$ permutations of the $\alpha$ are different from one another. If then $g(t) = 0$ is the Galois resolvent of $F(x) = 0$ whose root is this $\varrho$, then
\[
g(\varrho) = 0.
\]
